\noindent\textbf{Extended summary}\par\smallskip
\noindent Sports organisations test athletes on force plates routinely, but decisions usually rest on a few software-generated numbers such as jump height and peak force. This paper asks what those numbers leave out: at which points along the take-off and the landing of a countermovement jump do athletes differ, and do athletes with the same jump height move in the same way?

We analysed \nTrials{} jumps from \nAth{} male football players, most of them adolescents, tested during routine monitoring at a private sports physiotherapy and performance centre in India. Each phase was expressed as a body-weight-normalised force curve. Function-on-scalar regression, estimated in two independent ways with athlete-level uncertainty, located where force relates to jump height, age and body mass. Permutation tests on athlete-mean curves compared high and low jumpers and adults with under-18 players while controlling false positives over the whole curve. Matched-height comparisons, validated clustering and curve alignment followed.

Better jumpers produced more force in the final third of take-off: in the top versus bottom tertile the peak difference was \tertToD{} body weights, a standardised effect of \tertToG. Outside the last tenth of take-off the conventional variables explained little of the curve, and jumps of the same height were nearly as varied as the whole sample. Apparent age differences disappeared once body mass and jump height were taken into account. Landing differences were weaker in standardised terms and largely a matter of timing. Curve profiles formed a continuum rather than distinct strategies, yet \stabLd\% of athletes kept the same landing profile across sessions, against \chanceLd\% expected by chance.

For managers, the implications are to target the late push when training for height, report curve features beside single numbers, watch landing load in young players, adjust youth benchmarks for size and track landing habits. Any link to injury remains a hypothesis.
